% ─────────────────────────────────────────────────────────────────────────────
% Anexo I — Declaración de originalidad y autorización (ETS de Ingeniería – ICAI,
% Universidad Pontificia Comillas). Se coloca tras la portada y antes del Resumen.
% ─────────────────────────────────────────────────────────────────────────────
\thispagestyle{empty}
\begin{center}
  {\large\bfseries ANEXO I}\\[0.4em]
  {\large\bfseries Declaración de originalidad y autorización}\\[0.4em]
  {\normalsize Curso académico 2025-2026}
\end{center}

\vspace{1.2em}

\noindent\textbf{Declaración de originalidad}

\vspace{0.4em}

\noindent Declaro bajo mi responsabilidad que el Proyecto presentado con el título
\emph{\tituloTFM} en la ETS de Ingeniería – ICAI de la Universidad Pontificia Comillas
en el curso académico 2025-2026 es de mi autoría y no ha sido presentado con
anterioridad a otros efectos. El Proyecto no es plagio de otro, ni total ni
parcialmente, y la información que ha sido tomada de otros documentos está debidamente
referenciada.

\vspace{1.2em}

\noindent\textbf{Uso de Inteligencia Artificial}

\vspace{0.4em}

\noindent Declaro bajo mi responsabilidad que (indicar la opción correcta):

\vspace{0.6em}

\noindent$\square$\quad No he utilizado Inteligencia Artificial en la elaboración del
presente documento.

\vspace{0.5em}

\noindent$\boxtimes$\quad He utilizado Inteligencia Artificial en la elaboración del
presente documento y/o del Anexo B siempre en las condiciones permitidas por la
Universidad Pontificia Comillas, es decir, aplicando el Nivel~2 de la Escala de
Evaluación de Perkins et al. (2024): \enquote{La IA puede utilizarse para actividades
previas a la tarea, como la lluvia de ideas, la descripción y la investigación
inicial. Este nivel se centra en el uso de la IA para la planificación, las síntesis y
la generación de ideas, pero las evaluaciones deben hacer hincapié en la capacidad de
desarrollar y refinar estas ideas de forma independiente}. En concreto, la
Inteligencia Artificial ha sido empleada para:

\vspace{0.5em}

\noindent\fbox{\parbox{0.97\linewidth}{\vspace{0.2em}Apoyo en la revisión de estilo y
claridad de la redacción de la memoria. El diseño de los experimentos, el código, los
datos, los resultados y las conclusiones son trabajo del autor.\vspace{0.2em}}}

\vspace{1.6em}

\noindent\textbf{Autorización para la entrega del Proyecto}

\vspace{2.2em}

\noindent Firmado (alumno): \autorTFM\\[0.4em]
\noindent Fecha: 11.06.2026

\vspace{2.4em}

\noindent
\begin{minipage}[t]{0.48\linewidth}
  \centering
  El Director del Proyecto\\[2.6em]
  \rule{0.85\linewidth}{0.4pt}\\[0.3em]
  Fdo.: \tutorTFM\\[0.3em]
  Fecha:
\end{minipage}\hfill
\begin{minipage}[t]{0.48\linewidth}
  \centering
  El co-Director del Proyecto (si aplica)\\[2.6em]
  \rule{0.85\linewidth}{0.4pt}\\[0.3em]
  Fdo.:\\[0.3em]
  Fecha:
\end{minipage}

\cleardoublepage
